\documentclass{article}
\usepackage{graphicx} % Required for inserting images
\usepackage{caption} % Required for \captionof
\usepackage{amsmath}
\usepackage{float}
\usepackage{comment}
\usepackage[T1]{fontenc}
\usepackage[utf8]{inputenc}
\usepackage{lmodern}
\usepackage{microtype}
\usepackage{xcolor}
\usepackage{framed}
\renewenvironment{leftbar}{%
  \def\FrameCommand{\textcolor{orange}{\vrule width 3pt} \hspace{10pt}}%
  \MakeFramed {\advance\hsize-\width \FrameRestore}}%
 {\endMakeFramed}

\title{Mechanikai mérés}
\author{Kern Luca, Vincze Csongor}
\date{2026 Április 22}

\begin{document}

\maketitle

\section*{\documentclass{article}
\usepackage{graphicx} % Required for inserting images
\usepackage{caption} % Required for \captionof
\usepackage{amsmath}
\usepackage{float}
\usepackage{comment}
\usepackage[T1]{fontenc}
\usepackage[utf8]{inputenc}
\usepackage{lmodern}
\usepackage{microtype}
\usepackage{xcolor}
\usepackage{framed}
\renewenvironment{leftbar}{%
  \def\FrameCommand{\textcolor{orange}{\vrule width 3pt} \hspace{10pt}}%
  \MakeFramed {\advance\hsize-\width \FrameRestore}}%
 {\endMakeFramed}

\title{Mechanikai mérés}
\author{Kern Luca, Vincze Csongor}
\date{2026 Április 22}

\begin{document}

\maketitle

\section*{\documentclass{article}
\usepackage{graphicx} % Required for inserting images
\usepackage{caption} % Required for \captionof
\usepackage{amsmath}
\usepackage{float}
\usepackage{comment}
\usepackage[T1]{fontenc}
\usepackage[utf8]{inputenc}
\usepackage{lmodern}
\usepackage{microtype}
\usepackage{xcolor}
\usepackage{framed}
\renewenvironment{leftbar}{%
  \def\FrameCommand{\textcolor{orange}{\vrule width 3pt} \hspace{10pt}}%
  \MakeFramed {\advance\hsize-\width \FrameRestore}}%
 {\endMakeFramed}

\title{Mechanikai mérés}
\author{Kern Luca, Vincze Csongor}
\date{2026 Április 22}

\begin{document}

\maketitle

\section*{3. Feladat: Határozza meg a torziós asztal tehetetlenségi nyomatékát három különböző módon}

Mérőpár mindkét tagja mérje meg a paramétereket!

$R_{tárcs} = \text{m}$
$d_{tárcsa} = \text{m}$
$\rho_{alumínium} = 2700 \frac{kg}{m^3}$

\begin{table}[H]
    \centering
    \renewcommand{\arraystretch}{1.2}
    \begin{tabular}{|l|c|c|c|}
        \hline
        Mérőpár & $R_{tárcs}$ [m] & $d_{tárcsa}$ & $\Theta$ [kg m$^2$] \\ \hline
        Luca & & & \\ \hline
        Csongor & & & \\ \hline
    \end{tabular}
    \caption{Asztal tehetetlenségi nyomatékának meghatározása}
    \label{tab:hotan_2}
\end{table}


\renewcommand{\thesubsection}{\alph{subsection})} % Changes 1.1 to a)
\subsection{Számítással az asztal tömegének és sugarának ismeretében}

Mérőpár mindkét tagja mérje meg a paramétereket!

$R_{tárcs} = \text{m}$
$d_{tárcsa} = \text{m}$
$\rho_{alumínium} = 2700 \frac{kg}{m^3}$

\begin{table}[H]
    \centering
    \renewcommand{\arraystretch}{1.2}
    \begin{tabular}{|l|c|c|c|}
        \hline
        Mérőpár & $R_{tárcs}$ [m] & $d_{tárcsa}$ & $\Theta$ [kg m$^2$] \\ \hline
        Luca & & & \\ \hline
        Csongor & & & \\ \hline
    \end{tabular}
    \caption{Asztal tehetetlenségi nyomatékának meghatározása}
    \label{tab:hotan_2}
\end{table}


\subsection{A rugó direkciós nyomatékának, a lengésidőnek és a csillapítási tényezőnek az ismeretében}

\[
\theta = \left( \frac{T}{2\pi} \right)^2 \cdot k^*
\]


\subsection{Egy ismert tehetetlenségi nyomatékú tárcsa felhasználásával!}

\[
\theta = \theta_0 \frac{T^2}{T'^2 - T^2}.
\]

$m_{korong} = \text{kg}$
$r_{korong} = \text{m}$
$\theta_{korong} = \text{kg m}^2$
$m_{rézcsap} = 2,2 \text{kg}$


\end{document}